\documentclass{article}

\begin{document}

\title{Gaze2Guide: Memory-Augmented Gaze Reasoning for Exhibit-to-Exhibit Decision and Dwell-Time Prediction}
\author{Xijing Wang, Lei He, and Nianxiong Liu}
\date{}
\maketitle

\begin{abstract}
Museum guide and service robots need human goal and timing priors to decide where to guide and when to intervene without being intrusive. We propose IROS\_gaze, a memory-augmented gaze reasoning framework that predicts the visitor's next exhibit, dwell time, and attention level. We construct instruction-style training data from real eye-tracking collected in two venues (Tsinghua Art Museum and an Osaka exhibition), resulting in 8,159 high-quality labeled segments.
\end{abstract}

\section{INTRODUCTION}

Robots operating in public indoor environments must anticipate where a visitor will go next and how long they will stay. Our key idea is to combine an exhibit-level neighbor graph, short- and long-term memory, and LLM-based constrained reasoning to jointly predict next exhibit, dwell time, and attention level.

\textbf{Our contributions are:}

\begin{enumerate}
    \item A robot-ready gaze-driven prior framework
    \item A structured reasoning pipeline with dual-layer memory
    \item Robot-style evaluation through topological guide simulation
\end{enumerate}

\section{RELATED WORK}

Gaze has been widely used as a signal of human attention. Topological representations have proven effective for robot navigation. Our work combines gaze-driven attention with topological constraints for goal prediction.

\section{DATA AND PROBLEM SETUP}

We collected eye-tracking data in two venues: Tsinghua Art Museum (12 participants) and Osaka Exhibition (8 participants). Raw records (9,158) were cleaned to obtain 8,159 high-quality labeled segments.

\textbf{Attention Level Discretization:}

\begin{center}
\begin{tabular}{clc}
\hline
Level & Duration Range & Interpretation \\
\hline
A & 90--300s & Deep engagement \\
B & 45--90s & Sustained attention \\
C & 20--45s & Moderate interest \\
D & 8--20s & Brief interest \\
E & 3--8s & Glancing \\
\hline
\end{tabular}
\end{center}

\textbf{Dataset Statistics:}

\begin{center}
\begin{tabular}{lr}
\hline
Statistic & Value \\
\hline
Venues & Tsinghua Art Museum; Osaka Exhibition \\
Sampling Rate & 30 Hz \\
Raw / Clean Segments & 9,158 / 8,159 \\
Task Counts & T1: 3,038; T2: 3,038; T3: 2,083 \\
Unique Exhibits & 47 (TH: 28, OS: 19) \\
\hline
\end{tabular}
\end{center}

\section{METHOD}

IROS\_gaze consists of four layers: (1) Data Input Layer, (2) Memory Management Layer, (3) Reasoning \& Prediction Layer, and (4) Robot Interface Layer.

\subsection{Dual-Layer Memory}

\textbf{Short-Term Memory:} Maintains a sliding window of the most recent N=5 gaze records.

\textbf{Long-Term Memory:} Maintains global statistics including visit frequency, cumulative dwell time, and attention level distribution.

\subsection{Sequence Predictor}

The predictor performs K-step lookahead (K=5) by iteratively applying the single-step predictor to generate future viewing trajectories.

\section{EXPERIMENTS}

\subsection{Main Results}

\begin{center}
\begin{tabular}{lccccc}
\hline
Method & Hit@1 & MRR & Hit@3 & Dwell MAE & Attn Acc \\
\hline
Markov Chain & 45.2\% & 0.62 & 68.1\% & -- & -- \\
LSTM & 56.4\% & 0.71 & 78.2\% & 24.5s & 54.3\% \\
GPT-4o (Zero-Shot) & 58.3\% & 0.74 & 81.2\% & 18.3s & 61.2\% \\
\textbf{Ours (Full)} & \textbf{68.3\%} & \textbf{0.84} & \textbf{88.4\%} & \textbf{12.1s} & \textbf{72.4\%} \\
\hline
\end{tabular}
\end{center}

\subsection{Ablation Study}

\begin{center}
\begin{tabular}{lcccc}
\hline
Variant & Top-1 & Dwell MAE & Attn Acc & Regret \\
\hline
\textbf{Full (Ours)} & \textbf{68.3\%} & \textbf{12.1s} & \textbf{72.4\%} & \textbf{1.2} \\
-No Memory & 54.2\% & 18.3s & 61.2\% & 2.4 \\
-No Topology & 61.8\% & 15.7s & 68.1\% & 1.8 \\
\hline
\end{tabular}
\end{center}

Memory removal causes the largest drop, confirming the importance of temporal patterns.

\section{ROBOT SIMULATION}

We evaluate in a topological guide simulation with 19 exhibits and 37 spatial edges.

\begin{center}
\begin{tabular}{lcccc}
\hline
Policy & Hit@1 & Hit@3 & Timing Error & Regret \\
\hline
Random & 15.3\% & 42.1\% & 58.4s & 4.8 \\
Markov-Frequency & 45.2\% & 68.1\% & 32.1s & 2.6 \\
\textbf{Ours-Prior} & \textbf{68.3\%} & \textbf{88.4\%} & \textbf{14.7s} & \textbf{1.2} \\
\hline
\end{tabular}
\end{center}

Our prior achieves 23\% higher goal hit rate and halves timing error.

\section{CONCLUSION}

We presented IROS\_gaze, a memory-augmented gaze reasoning framework for museum guide robots. Our system achieves strong performance on next-step prediction and demonstrates utility in robot-facing simulation.

\end{document}
